% paper/sections/10_implementation.tex
\section{演算子整合性の検証：Balanced-Force と DCCD-GFM 実装の確認}
\label{sec:impl_collocate}

第~\ref{sec:algorithm}章（完全アルゴリズムフロー）で示した時間積分ループでは，
離散演算子の整合性が安定動作と高精度の前提となる．
第~\ref{sec:collocate}章（§\ref{sec:rhie_chow}，§\ref{sec:balanced_force}）は
Balanced-Force の理論的根拠と必要条件を導出した．
本章はその\textbf{実装が正しく機能しているかを検証する}ことを目的とし，
静止液滴テストによる定量的確認手順を示す．

% ─────────────────────────────────────────────────────────────────────────────
\subsection{RC-CCD 関係の明確化：安定性と高精度の役割分担}
\label{sec:impl_bf_consistency}
% ─────────────────────────────────────────────────────────────────────────────

第~\ref{sec:collocate}章の結論を実装の観点から整理する．

\textbf{チェッカーボード抑制（安定性）：} 本稿では DCCD 散逸機構（§\ref{sec:dccd_decoupling}，$\varepsilon_d=1/4$）が
コロケート格子固有のチェッカーボード不安定性を抑制する．
Rhie--Chow 補正（付録~\ref{app:rhie_chow_full}）は参照実装として記録する．

\textbf{Balanced-Force 精度：} 寄生流れの抑制は離散演算子の整合性による．
§\ref{sec:balanced_force} の議論より，離散演算子不整合（式~\eqref{eq:bf_operator_mismatch}）をゼロにするには，
圧力勾配・速度補正の 2 箇所に同一の $D^{(1)}_\mathrm{CCD}$ を用いなければならない：

\begin{enumerate}
  \item 予測ステップの IPC 項（前時刻圧力勾配）：
        $-D^{(1)}_\mathrm{CCD} p^n$
  \item 速度補正ステップ：
        $-({\Delta t}/{\rho})\,D^{(1)}_\mathrm{CCD}(\delta p)$
\end{enumerate}

表面張力は GFM により PPE 右辺に直接組み込まれる（§\ref{sec:gfm}）ため，
予測ステップに体積力項は存在しない．
2箇所の演算子が統一されることで，式~\eqref{eq:bf_operator_mismatch} の不整合はゼロに設計されている．

\textbf{CCD 曲率精度が Balanced-Force の上限を決める：}
寄生流れの下限は曲率計算の精度によって決まる．
CCD $\Ord{h^6}$ による曲率評価と GFM 鋭利界面処理（§\ref{sec:gfm}）の組み合わせにより，
CSF モデル誤差（$\Ord{h^2}$）を根本的に除去し，より低い寄生流れが期待できる
（定量評価は §\ref{sec:impl_bf_verify} および §\ref{sec:bench_droplet}）．

% ─────────────────────────────────────────────────────────────────────────────
\subsection{Balanced-Force 条件の定量評価：監視すべき指標}
\label{sec:impl_bf_rc}
% ─────────────────────────────────────────────────────────────────────────────

実装の正しさを確認するため，以下の指標を監視する：

\begin{enumerate}
  \item \textbf{$O(h^q)$ 収束次数の確認：}
        格子幅 $h$ を変えながら最大速度振幅 $U_{\mathrm{max}}(h) = \|\bm{u}\|_{L^\infty}$ を測定し，
        $U_{\mathrm{max}} \to 0$ の収束を確認する．
        GFM + DCCD の組み合わせでは CSF モデル誤差（$\Ord{h^2}$）が除去されるため，
        より高次の収束（$q > 2$）が期待される（§\ref{sec:bench_droplet} 参照）．
  \item \textbf{演算子整合性チェック（コード検証）：}
        Predictor の $-D^{(1)}_\mathrm{CCD} p^n$ と Corrector の $-D^{(1)}_\mathrm{CCD}(\delta p)$
        が同一の差分演算子インスタンスを使っているかコードレベルで確認する．
  \item \textbf{GFM 補正項の符号・スケール確認：}
        界面を挟む格子点ペアで $b_i^{\mathrm{GFM}}$（式~\eqref{eq:gfm_rhs_correction}）の
        符号と大きさが理論値（$\kappa_f/We \cdot (1/\rho)^{\mathrm{harm}}/h^2$）と一致するか確認する．
\end{enumerate}

% ─────────────────────────────────────────────────────────────────────────────
\subsection{Balanced-Force 実装の確認方法：静止液滴テスト}
\label{sec:impl_bf_verify}
% ─────────────────────────────────────────────────────────────────────────────

\begin{tcolorbox}[algbox, title={静止液滴テスト：Balanced-Force の検証手順}]
\begin{enumerate}
  \item \textbf{初期条件：} 計算領域中央に半径 $R$ の円形液滴を置く．
        速度 $\bm{u} = \bm{0}$，圧力は Young-Laplace 解
        $p^0 = -\sigma\kappa_0 = \sigma/R$（液相内，$\kappa_0 = -1/R$），
        $p^0 = 0$（気相）．

  \item \textbf{時間発展：} 数十タイムステップ計算する．
        Balanced-Force が正しく実装されていれば，速度は CSF モデル誤差が律速する
        $O(\varepsilon^2) \approx O(h^2)$ レベルに留まる
        （例：$N=64$ で $\|\bm{u}\|_{L^\infty} \lesssim 10^{-3}$〜$10^{-4}$ 程度）．
        CSF が界面を有限厚さ $\varepsilon \sim h$ で近似する限り，
        \textbf{モデル誤差 $O(h^2)$} が寄生流れの下限を決めるため，
        機械精度への収束は物理的に不可能である（付録~\ref{app:balanced_force_taylor}）．

  \item \textbf{収束確認：} 格子幅 $h$ を変えながら最大速度振幅
        \[
          U_{\mathrm{max}}(h) = \|\bm{u}\|_{L^\infty}
        \]
        を測定し，$U_{\mathrm{max}}(h) = O(h^q)$（$q \approx 2$）の収束を確認する．
        CSF モデルの界面厚さ $\varepsilon \sim O(h)$ に起因するモデル誤差 $O(\varepsilon^2) \approx O(h^2)$
        が律速するため，Balanced-Force + CCD を正しく実装しても収束次数は $q \approx 2$ となる
        （付録~\ref{app:balanced_force_taylor}，§3）．
        Balanced-Force の効果は収束\textbf{次数}の改善ではなく，収束\textbf{係数}の大幅な低減
        ——例えば $U_{\mathrm{max}}$ が $O(10^{-1})$〜$O(1)$（実装ミス）から
        $O(10^{-3})$〜$O(10^{-4})$（正しい実装）へ——として現れる（§\ref{sec:bench_droplet} 参照）．

  \item \textbf{失敗の兆候：} $U_{\mathrm{max}}$ が $O(1)$〜$O(10^{-2})$ 程度で収束しない場合，
        以下を確認する：
        \begin{itemize}
          \item 圧力勾配と表面張力に同一の差分演算子を使っているか？
          \item Rhie--Chow の密度係数と PPE の係数に同じ調和平均を使っているか？
          \item 曲率計算と表面張力の計算タイミングが整合しているか？
          \item 表面張力支配の問題（$We \ll 1$）では，式~\eqref{eq:rc-face-balanced} の
                表面張力補正項を Rhie--Chow に追加しているか（§\ref{sec:rc_balanced_force} 参照）？
        \end{itemize}
\end{enumerate}
\end{tcolorbox}
